%! Confidence Intervals for the Difference of Two Means — Unit 7, Lesson 7.6 — Exit Ticket (Answer Key)
\documentclass[10pt]{article}
\usepackage{apstats-article}
\usepackage{apstats-key}   % pulls in -boxes; do NOT also load -boxes
\newcommand{\TallMath}[1]{$\displaystyle #1\rule[-1.4em]{0pt}{3.2em}$}

\begin{document}

\pageheader{Unit 7, Lesson 7.6}{Exit Ticket --- Answer Key}
\namedateperiod

\vspace{0.15in}

A nutritionist wants to estimate the difference in mean daily sugar intake (grams) between
adults in two cities. City X: $n_1 = 35$, $\bar x_1 = 82$ g, $s_1 = 14$ g. \quad
City Y: $n_2 = 40$, $\bar x_2 = 75$ g, $s_2 = 11$ g.
Both samples are random; both sample distributions are roughly symmetric with no outliers.
Technology gives $df \approx 65$, $t^* \approx 1.997$ (95\% confidence).

\begin{enumerate}
  \item Identify the procedure and write the point estimate.

        Procedure: \ans{two-sample $t$-interval for $\mu_{\text{X}} - \mu_{\text{Y}}$}

        Point estimate: $\bar x_1 - \bar x_2 = $ \ans{$82 - 75 = 7$ g}

  \item Verify both conditions.
        \begin{itemize}
          \item Random / Independence: \ansline{Random samples from two independent groups --- met.}
          \item Normal: \ansline{$n_1 = 35 \ge 30$ and $n_2 = 40 \ge 30$; CLT applies for both --- met.}
        \end{itemize}

  \item Compute the standard error.

        $SE = \sqrt{\dfrac{s_1^2}{n_1} + \dfrac{s_2^2}{n_2}}
            = \sqrt{\dfrac{({\color{keyred}\mathbf{14}})^2}{{\color{keyred}\mathbf{35}}} + \dfrac{({\color{keyred}\mathbf{11}})^2}{{\color{keyred}\mathbf{40}}}} \approx$ \ans{2.939}

  \item Construct and interpret the 95\% confidence interval.

        $CI = {\color{keyred}\mathbf{7}} \pm 1.997 \times {\color{keyred}\mathbf{2.939}} = ($ \ans{1.13} , \ans{12.87} $)$

        \textit{Interpretation:} \ansline{We are 95\% confident the true difference in mean daily sugar intake (City X $-$ City Y) is between 1.13 g and 12.87 g.}
        \ansline{Since the interval is entirely positive and does not contain 0, this provides convincing evidence that adults in City X consume more sugar on average than adults in City Y.}

        \begin{teachernote}
        $SE = \sqrt{196/35 + 121/40} = \sqrt{5.600 + 3.025} = \sqrt{8.625} \approx 2.937$.
        $ME = 1.997 \times 2.937 \approx 5.865$. CI: $(7 - 5.87,\; 7 + 5.87) = (1.13,\; 12.87)$.
        \end{teachernote}
\end{enumerate}

\end{document}
